\documentclass[11pt]{article}

\usepackage[utf8]{inputenc}
\usepackage[T1]{fontenc}
\usepackage{geometry}
\usepackage{amsmath}
\usepackage{amssymb}
\usepackage{booktabs}
\usepackage{hyperref}
\usepackage{xurl}
\usepackage{xcolor}
\usepackage{listings}
\usepackage{enumitem}
\usepackage{parskip}
\usepackage{graphicx}
\graphicspath{{../images/}}

\geometry{margin=1in}

\lstset{
  language=Java,
  basicstyle=\ttfamily\small,
  keywordstyle=\color{blue},
  commentstyle=\color{gray},
  stringstyle=\color{teal},
  showstringspaces=false,
  columns=fullflexible,
  frame=single,
  framerule=0.2pt,
  breaklines=true
}

\setlist[itemize]{topsep=0.3em,itemsep=0.2em}
\setlist[enumerate]{topsep=0.3em,itemsep=0.2em}

% Simple per-section source line.
\newcommand{\notesources}[1]{\par\noindent{\small\color{gray}\textit{Sources:} \sloppy #1\par}}

% Unit-wise source strings for this subject (used by slides/notes).
% Path is relative to the lecture `latex/` folder (the compilation CWD).
% OOPS with Java (CSEG1044) — unit-wise sources
%
% These are short, student-facing reference strings intended to be used with:
%   - \slidesources{...}  (in beamer slides)
%   - \notesources{...}   (in lecture notes)
%
% Style inspiration: other subjects in My-Subjects/ use semicolon-separated sources
% and include an "accessed" date for web documentation.

\newcommand{\OOPSUnitISources}{Schildt, \textit{Java: A Beginner's Guide} (10e), McGraw-Hill (Java fundamentals + OOP basics).; Oracle Java Tutorials: Getting Started + Language Basics + Classes and Objects (accessed \today).; Java SE API docs (e.g., \texttt{String}, \texttt{Scanner}) (accessed \today).}

\newcommand{\OOPSUnitIISources}{Schildt, \textit{Java: The Complete Reference} (12e), McGraw-Hill (classes, inheritance, packages, interfaces).; Bloch, \textit{Effective Java} (3e), Addison-Wesley, 2018 (design choices; interfaces vs abstract classes).; Oracle Java Tutorials: Classes and Objects + Interfaces + Packages (accessed \today).; Java SE API docs (e.g., \texttt{Comparable}, \texttt{Runnable}) (accessed \today).}

\newcommand{\OOPSUnitIIISources}{Schildt, \textit{Java: The Complete Reference} (12e) (nested/inner classes, exceptions, threads, I/O).; Oracle Java Tutorials: Nested Classes + Exceptions + Concurrency + Basic I/O (accessed \today).; Java SE API docs (e.g., \texttt{Thread}, \texttt{AutoCloseable}) (accessed \today).}

\newcommand{\OOPSUnitIVSources}{Schildt, \textit{Java: The Complete Reference} (12e) (generics, lambdas, Swing, JDBC).; Bloch, \textit{Effective Java} (3e) (generics + lambdas).; Oracle Java Tutorials: Generics + Lambda Expressions + Swing + JDBC Basics (accessed \today).; Java SE API docs (e.g., \texttt{java.util.function}, \texttt{javax.swing}, \texttt{java.sql}) (accessed \today).}

\newcommand{\OOPSUnitVSources}{Schildt, \textit{Java: The Complete Reference} (12e) (Collections Framework + wrapper classes).; Oracle Java docs: Collections Framework overview (accessed \today).; Java SE API docs (\texttt{java.util} collections; wrapper classes in \texttt{java.lang}) (accessed \today).}

\newcommand{\OOPSUnitVISources}{Git SCM: \textit{Pro Git} (2e) (version control workflow), accessed \today.; GitHub Docs (repositories, issues, pull requests), accessed \today.; Oracle Java Tutorials: Swing + JDBC (accessed \today).; JUnit 5 User Guide (accessed \today).; Apache Maven Project: Getting Started (accessed \today).; Gradle User Manual: Getting Started (accessed \today).; UML class diagrams: UML 2.x overview/spec (accessed \today).}

\newcommand{\OOPSMiscWhyInterfacesSources}{Schildt, \textit{Java: The Complete Reference} (12e) (interfaces).; Bloch, \textit{Effective Java} (3e) (prefer interfaces where appropriate).; Oracle Java Tutorials: Interfaces (accessed \today).; Java SE API docs (e.g., \texttt{Comparable}, \texttt{Runnable}) (accessed \today).}



\title{Object Oriented Programming (CSEG1044)\\Unit 04 -- Lecture 03 Notes\\Lambdas + Functional Interfaces}
\begin{document}
\maketitle
\tableofcontents

\section{Lecture Snapshot}
\begin{itemize}
  \item \textbf{Unit focus:} Generics, Lambdas, Swing, and JDBC
  \item \textbf{Main new ideas:} Functional interfaces (concept), Target typing (how lambdas get a type), \texorpdfstring{Anonymous class $\rightarrow$ lambda}{Anonymous class -> lambda}
  \item \textbf{Course thread:} Use Lambdas + Functional Interfaces to move Campus Resource Manager toward a real desktop application with UI, callbacks, and persistence.
\end{itemize}
\notesources{\OOPSUnitIVSources}

\section{Lecture Overview}
Lambdas are a short way to write ``behavior as a value'' in Java.
They work with \textbf{functional interfaces} (interfaces with exactly one abstract method).
This is useful for:
\begin{itemize}
  \item sorting (\texttt{Comparator}),
  \item callbacks and event handling (Swing listeners),
  \item threading (\texttt{Runnable}),
  \item cleaner APIs for reusable behavior.
\end{itemize}
\notesources{\OOPSUnitIVSources}

\section{Prerequisite Bridge}
Before reading the new material in depth, do a fast recall of the older ideas listed below.
\begin{itemize}
  \item \textbf{Interfaces}: A lambda only works when a target functional interface defines exactly one abstract method. Main reference: \texttt{Unit 02 / Lecture 05 slides/notes}.
  \item \textbf{Anonymous classes}: Lambdas are easier to appreciate once a more verbose callback style is already known. Main reference: \texttt{Unit 03 / Lecture 01 slides/notes}.
\end{itemize}
This lecture only revisits those ideas briefly so that the main explanation can stay focused on the new concept sequence.
\notesources{\OOPSUnitIVSources}

\section{Learning Outcomes}
After this lecture, you should be able to:
\begin{itemize}
  \item Identify functional interfaces and explain target typing.
  \item Write lambdas for common patterns (comparators, callbacks, runnable tasks).
  \item Apply basic readability rules for when to use lambdas.
\end{itemize}

\section{Suggested Reading Order}
\begin{enumerate}
  \item Read the \textbf{Prerequisite Bridge} first and confirm each recalled idea.
  \item Study the central explanation in this order: \textit{Functional interfaces (concept)} $\rightarrow$ \textit{Target typing (how lambdas get a type)} $\rightarrow$ \textit{\texorpdfstring{Anonymous class $\rightarrow$ lambda}{Anonymous class -> lambda}} $\rightarrow$ \textit{Lambda syntax (common forms)} $\rightarrow$ \textit{Common patterns}.
  \item Trace the worked example line by line before checking short answers.
  \item Attempt the practice section without looking at the solution immediately.
  \item Finish with the exit questions and further-study suggestions to confirm retention.
\end{enumerate}
\notesources{\OOPSUnitIVSources}

\section{Key Terms (Quick)}
\begin{itemize}
  \item \textbf{Functional interface}: interface with 1 abstract method (SAM).
  \item \textbf{SAM}: Single Abstract Method.
  \item \textbf{Lambda}: expression like \texttt{(a, b) -> a + b}.
  \item \textbf{Target typing}: compiler infers lambda type from context (assignment/argument).
  \item \textbf{Effectively final}: local variable not modified after initialization (can be captured in lambda).
\end{itemize}

\section{Detailed Notes}
\subsection{Functional interfaces (concept)}
A functional interface has exactly one abstract method.
It can still have:
\begin{itemize}
  \item \texttt{default} methods,
  \item \texttt{static} methods.
\end{itemize}

\textbf{Examples:}
\begin{itemize}
  \item \texttt{Runnable} (method \texttt{run()})
  \item \texttt{Comparator<T>} (method \texttt{compare(a,b)})
  \item \texttt{ActionListener} in Swing (method \texttt{actionPerformed(e)})
\end{itemize}

\paragraph{Annotation (optional but useful).}
\texttt{@FunctionalInterface} tells the compiler to enforce the SAM rule.

\subsection{Target typing (how lambdas get a type)}
Lambda expressions do not have a type by themselves.
They get a type from where they are used:
\begin{itemize}
  \item assignment: \texttt{Runnable r = () -> ...}
  \item method argument: \texttt{button.addActionListener(e -> ...)}
\end{itemize}
If the context expects a functional interface, the lambda is converted to that interface.

\subsection{\texorpdfstring{Anonymous class $\rightarrow$ lambda}{Anonymous class -> lambda}}
Think of a lambda as the short form of a one-method anonymous class.
\begin{itemize}
  \item Anonymous class: write the interface name, the method header, and the body.
  \item Lambda: keep only the parameters and the body because the functional interface already gives the method name.
  \item Use this shortcut only when the behavior is small enough to stay readable.
\end{itemize}

\subsection{Lambda syntax (common forms)}
\begin{itemize}
  \item No parameters: \texttt{() -> System.out.println("Hi")}
  \item One parameter: \texttt{x -> x * 2} (parentheses optional)
  \item Many parameters: \texttt{(a, b) -> a + b}
  \item Block body: \texttt{x -> \{ int y = x + 1; return y; \}}
\end{itemize}

\subsection{Common patterns}
\subsubsection{Comparator for sorting}
Instead of writing a full comparator class, you can write:
\begin{itemize}
  \item \texttt{(a, b) -> a.getCgpa() < b.getCgpa() ? 1 : -1}
\end{itemize}
or use helper APIs like \texttt{Comparator.comparing}.

\subsubsection{Runnable for tasks/threads}
You can write a runnable task in one line and pass it to a \texttt{Thread}.

\subsubsection{Callbacks/listeners}
Swing event handling uses listeners, which are functional interfaces in many cases (or have adapter patterns).
Lambdas make event handling code shorter.

\subsection{Readability guidelines}
Use lambdas when:
\begin{itemize}
  \item the behavior is short (1--3 lines),
  \item the logic is clear from context,
  \item you do not need to reuse it elsewhere.
\end{itemize}
Prefer a named method/class when:
\begin{itemize}
  \item the lambda becomes long/complex,
  \item you need to unit-test or reuse the behavior,
  \item error handling and branching make it hard to read.
\end{itemize}
\notesources{\OOPSUnitIVSources}

\section{Running Project Connection}
Use Lambdas + Functional Interfaces to move Campus Resource Manager toward a real desktop application with UI, callbacks, and persistence.
\begin{itemize}
  \item Use Campus Resource Manager as the running case study, or map the same idea to your own project domain if you are using another example.
  \item Ask where today's idea should live: model, service, DAO, UI, or team workflow.
  \item Use the lecture example as a smaller version of the same design choice you will make in the capstone.
\end{itemize}
\notesources{\OOPSUnitIVSources}

\section{Common Mistakes and Confusions}
\begin{itemize}
  \item Assuming any interface can be used as a lambda target.
  \item Packing too much business logic into one lambda.
  \item Ignoring the target type and getting confused about the method being implemented.
\end{itemize}
\notesources{\OOPSUnitIVSources}

\section{Worked Example: sorting with lambda}
\begin{lstlisting}
import java.util.ArrayList;
import java.util.Comparator;
import java.util.List;

class Student {
    final String rollNo;
    final double cgpa;
    Student(String rollNo, double cgpa) { this.rollNo = rollNo; this.cgpa = cgpa; }
    public double getCgpa() { return cgpa; }
    public String toString() { return rollNo + ":" + cgpa; }
}

public class Demo {
    public static void main(String[] args) {
        List<Student> list = new ArrayList<>();
        list.add(new Student("21CS1001", 8.2));
        list.add(new Student("21CS1002", 9.1));
        list.add(new Student("21CS1003", 7.6));

        list.sort((a, b) -> Double.compare(b.getCgpa(), a.getCgpa())); // desc
        System.out.println(list);

        list.sort((a, b) -> Double.compare(a.getCgpa(), b.getCgpa())); // asc
        System.out.println(list);
    }
}
\end{lstlisting}

\section{Demo / Observation Task}
Use this block as a short reproducible demo or observation task.
\begin{itemize}
  \item Reproduce the lecture's main example around \textit{Functional interfaces (concept)} once without looking at the solution first.
  \item Change one input, rule, or design choice in Campus Resource Manager and predict the result before checking it.
  \item Record one success case, one edge case, and one failure case so the idea becomes observable instead of purely theoretical.
\end{itemize}
\notesources{\OOPSUnitIVSources}

\section{Practice Check (with brief solutions)}
\begin{enumerate}
  \item What is a functional interface?\par
  \textbf{Answer:} an interface with exactly one abstract method (SAM).

  \item Give one example of where you can use lambdas.\par
  \textbf{Answer:} sorting with \texttt{Comparator}, or creating a \texttt{Runnable}.

  \item Can a lambda capture a local variable that changes later?\par
  \textbf{Answer:} No, captured local variables must be effectively final.
\end{enumerate}

\section{Self-Study Practice (no solutions)}
\begin{enumerate}
  \item Write a lambda for \texttt{Runnable} that prints numbers 1--5.
  \item Sort a list of strings by length using a lambda comparator.
  \item Replace a complex lambda in your code with a named method and explain why it is clearer.
\end{enumerate}

\section{Further Study}
\begin{itemize}
  \item Revisit the prerequisite references if any recall bullet still feels weak.
  \item Rework the lecture example with one changed assumption, input, or design choice.
  \item Read the textbook and documentation references listed in the source line for a second explanation of the same topic.
  \item Apply the same idea once inside Campus Resource Manager so that the concept moves from theory to design practice.
\end{itemize}
\notesources{\OOPSUnitIVSources}

\section{Exit Questions (with sample answers)}
\textbf{Q1.} What is target typing?\par
\textbf{Sample answer:} Target typing is how the compiler infers a lambda's type from the expected functional interface in the context (assignment or method parameter).

\vspace{0.6em}
\textbf{Q2.} When should you avoid using a lambda?\par
\textbf{Sample answer:} Avoid lambdas when the logic becomes long or complex, because it hurts readability and reuse. Prefer a named method/class in that case.

\vspace{0.6em}
\textbf{Q3.} Why must a functional interface contain exactly one abstract method?\par
\textbf{Sample answer:} Because the lambda expression needs one unambiguous behavior contract to implement.

\end{document}
